% Generated draft: requires mathematical and source review.
\documentclass[11pt]{article}
\usepackage[margin=25mm]{geometry}
\usepackage{fontspec}
\setmainfont{Latin Modern Roman}
\usepackage{amsmath,amssymb,mathtools,unicode-math}
\setmathfont{Latin Modern Math}
\usepackage{enumitem,needspace,xurl,hyperref}
\setlength{\parindent}{0pt}
\setlength{\parskip}{5pt}
\setlength{\emergencystretch}{3em}
\newcommand{\N}{\mathbb{N}}
\newcommand{\Z}{\mathbb{Z}}
\newcommand{\Q}{\mathbb{Q}}
\newcommand{\R}{\mathbb{R}}
\newcommand{\C}{\mathbb{C}}
\newcommand{\F}{\mathbb{F}}
\newcommand{\A}{\mathbb{A}}
\newcommand{\PP}{\mathbb{P}}
\newcommand{\E}{\mathbb{E}}
\newcommand{\eps}{\varepsilon}
\newcommand{\dd}{\,\mathrm{d}}
\begin{document}
\section*{UnsolvedMath ID: 20000448}
\textbf{Source problem code:} catalogue entry
\subsection{Definitions and mathematical statement}
For a field $F$ of characteristic zero, let
\[
 V_n(F)=\left\{\sum_{i=1}^n\frac{2}{x_i^2-2}:x_i\in F,\ x_i^2\ne2\right\},\qquad n\geq1.
\]
Poonen's AIM question asks to characterize $V_n(\Q)$ for fixed $n$. It states the necessary integrality conditions $|\alpha|_p\leq1$ at $p=2$ and at every odd prime with Legendre symbol $(2/p)=-1$, and asks whether these suffice. The formula is reconstructed from damaged fraction/superscript formatting; the source explicitly takes finite $x_i\in\Q$, not points at infinity.

The report claims for every odd prime with $(2/p)=-1$ that
\[
 V_2(\Q_p)=\Z_p.
\]
At the ramified prime it gives only the following exact residue images:
\[
 V_1(\Q_2)\bmod8=\{0,1,6,7\},\qquad
 V_2(\Q_2)\bmod8=\{0,1,2,4,5,6,7\}.
\]
Thus $\alpha=3$ satisfies all printed integrality conditions but is not in $V_2(\Q_2)$ or $V_2(\Q)$. In contrast $3\in V_3(\Q)$, for instance all three variables can equal $2$. This refutes bare integrality sufficiency for $n=2$, not the Hasse principle, since the counterexample already fails local solubility. The residue assertion does not classify every $2$-adic value in each allowed residue class.

The report also gives an exact one-summand criterion, which in this convention is
\[
 \alpha\in V_1(\Q)\quad\Longleftrightarrow\quad
 \alpha\ne0\ \text{and}\ 2+2/\alpha\in\Q^2.
\]
Its norm-one-torus reformulation uses $L=\Q(\sqrt2)$ and $u=(x+\sqrt2)/(x-\sqrt2)$, with $N_{L/\Q}(u)=1$ and $u\ne1$. The trace identity yields
\[
 \alpha\in V_n(\Q)\quad\Longleftrightarrow\quad
 \exists u_1,\ldots,u_n\in L^\times:\ N(u_i)=1,\ u_i\ne1,
 \quad\sum_i\operatorname{Tr}_{L/\Q}(u_i)=4\alpha+2n.
\]
Zero-sum padding gives the reported length inclusions; explicitly $1+(-1)=0$ and $1+1+(-2)=0$ are values at $(2,0)$ and $(2,2,1)$, so $V_n\subset V_{n+2}$ and $V_n\subset V_{n+3}$. No general fixed-length global classification follows from these reductions.

The source remarks suggest analogous Waring-type questions for rational functions without rational poles, note unspecified earlier work when poles are rational, and propose the local-global formulation (Swinnerton-Dyer). They suggest local solubility for sufficiently large fixed $n$ and a Diophantine-definition application if the full integrality characterization holds: integrality at the inert half of the primes, together with $2$. Those are motivations, not established consequences of the two-summand result.
\subsection{Short English statement}
Characterize fixed-length sums of $2/(x^2-2)$ over $\Q$. The report shows that the stated integrality conditions miss a genuine $2$-adic obstruction for two summands: $3$ is excluded, although three summands represent it. It also supplies a one-term test, a norm-torus trace reduction and padding identities, leaving the full local-global question separate.
\subsection{Sources}
\begin{itemize}
\item Local registry: \url{data/unsolvedmath/problems.json}, permanent ID $20000448$, source code \texttt{catalogue entry}; original question, every remark and supplied research summary.
\item AIM workshop ``Rational and integral points on higher dimensional varieties'', Question 15, attributed to Poonen: \url{https://aimath.org/WWN/qptsurface2/qptsurface2.pdf}. Canonical location \texttt{aim-algebraic-number-theory-notes.json}, \texttt{notes[99]}; original literature field empty. No PDF was opened here.
\item Original remarks attribute the Hasse-principle reformulation to Swinnerton-Dyer and the rational-function/Diophantine motivations to Poonen; no references for the mentioned previous work are supplied.
\item Supplied machine-generated research attempt 1: partial result, problem open at the run and currently partially solved; rigor metadata ``proved'', novelty confidence low. The non-peer-reviewed local and torus results require independent verification before citation.
\end{itemize}
\end{document}
